\section{Conclusion}
In this C-MAPSS stress test, three bounded conclusions hold. First, coverage and the share of out-of-range inputs do not tell which way an interval-based maintenance rule fails: at matched excursion, drift and bias differed widely in coverage, and a sign-coherent positive shift gave premature triggers in \JDriftPremRange\% of evaluated six-condition engines at a 20-cycle lead time with no unrecognised need, while the opposite sign gave unrecognised needs. Second, output swaps identified the mean shift as the main driver of these errors, with scale changes shifting premature rates by between \JSigExtMinPt{} and \JSigExtMaxPt{} points under drift; under additive noise the scale effect has no order-independent sign. Third, no interval mechanism is recommendable on cost alone, since cost rankings changed across cost ratios in \JRevTotal{} of \JRevGroups{} settings and not at all under drift. Evaluations of intervals used for maintenance should therefore report both trigger errors with engine-level uncertainty; whether the same holds on recorded flight data or other rules remains open.
